\documentclass[a4paper,12pt]{article}
\usepackage[T2A]{fontenc}
\usepackage[utf8]{inputenc}
\usepackage[russian]{babel}
\usepackage{amsmath,amssymb,mathtools}
\usepackage{helvet}
\renewcommand{\familydefault}{\sfdefault}
\usepackage{geometry}
\geometry{left=23mm,right=22mm,top=24mm,bottom=24mm}
\usepackage{microtype}
\usepackage{tikz}
\usepackage{booktabs,tabularx,array}
\usepackage{enumitem}
\usepackage{parskip}
\usepackage{fancyhdr}
\usepackage{xcolor}
\usepackage[colorlinks=true,linkcolor=teal!65!black,bookmarks=false]{hyperref}
\definecolor{accent}{HTML}{087F87}
\definecolor{darkaccent}{HTML}{124F55}
\definecolor{errorred}{HTML}{9E322F}
\definecolor{quiet}{HTML}{5D6568}
\definecolor{maintext}{HTML}{242A2C}
\color{maintext}
\setlength{\parindent}{0pt}
\setlength{\parskip}{6pt}
\setlist[itemize]{leftmargin=*,itemsep=3pt,topsep=3pt}
\renewcommand{\arraystretch}{1.35}
\setlength{\headheight}{14pt}
\pagestyle{fancy}
\fancyhf{}
\fancyhead[L]{\small\color{quiet} МАТЕМАТИКА \hspace{0.7em} / \hspace{0.7em} РЕВЬЮ ДОМАШНЕЙ РАБОТЫ}
\fancyhead[R]{\small\color{quiet} 08.11.2026}
\fancyfoot[C]{\scriptsize\color{quiet} Лёвин Артём Александрович эксклюзивно для Николь Саркисьянц}
\fancyfoot[R]{\small\color{quiet}\thepage}
\renewcommand{\headrulewidth}{0.3pt}
\renewcommand{\footrulewidth}{0pt}
\newcommand{\problem}[2]{\phantomsection\label{task:#1}{\Large\color{darkaccent}\bfseries Задание № #1}\par\vspace{3pt}{\color{quiet}\small #2}\par\vspace{8pt}}
\newcommand{\blockhead}[1]{\vspace{8pt}{\bfseries\color{darkaccent}#1}\par\vspace{1pt}}
\newcommand{\err}[1]{\textcolor{errorred}{\textbf{Ошибка:} #1}}
\newcommand{\idea}[1]{\textcolor{darkaccent}{\textbf{Ключевая идея:} #1}}
\newcommand{\result}[1]{\textcolor{accent}{\textbf{Итог:} #1}}
\newcommand{\practice}[1]{\vspace{9pt}\par\noindent\textbf{\color{darkaccent}Задача на закрепление}\par #1}
\begin{document}
\thispagestyle{empty}
\begin{center}
\vspace*{25mm}
{\small\color{darkaccent}\MakeUppercase{Математика}}\par
\vspace{22mm}
{\LARGE Ревью домашней работы}\par
\vspace{13mm}
{\Large\color{accent} Задачи профильной математики}\par
\vspace{5mm}
{\normalsize\color{quiet} Ошибки, подробные решения и закрепление}\par
\vspace{43mm}
{\large Николь Саркисьянц}\par
\vspace{12mm}
{8 ноября 2026 г.}\par
\vfill
{\small Лёвин Артём Александрович}\par
{\small эксклюзивно для Николь Саркисьянц}\par
\vspace{19mm}
\begin{tikzpicture}
\draw[accent,line width=0.8pt] (0,0) .. controls (1.2,0.48) and (2.2,-0.40) .. (3.4,0)
.. controls (4.5,0.36) and (5.4,-0.42) .. (6.6,0);
\end{tikzpicture}
\end{center}
\clearpage
\section*{Сводная таблица}
В таблице приведены только задания, требующие исправления или дополнения. Успешно выполненные задания в документ не включены. Зачёркнутые на бланке задания рассматриваются как исключённые из варианта.\par
\vspace{8pt}
\small
\begin{tabularx}{\linewidth}{@{}p{13mm}p{24mm}X p{25mm}p{29mm}@{}}
\toprule
\textbf{№} & \textbf{Тема} & \textbf{Суть затруднения} & \textbf{Ваш ответ} & \textbf{Верный ответ}\\
\midrule
\hyperref[task:6]{6} & Вероятность & Сумма выплат принята за матожидание & 900\,000 руб. & 18 руб.\\
\hyperref[task:10]{10} & Прикладная задача & Ошибка при вычислении квадрата числа & 60,15 м & 60 м\\
\hyperref[task:11]{11} & Совместная работа & Решение и ответ отсутствуют & Не указан & 4 ч\\
\hyperref[task:16]{16} & Логарифмы & Ошибка после преобразования и пропущена ОДЗ & Отрезок & Два промежутка\\
\hyperref[task:17]{17} & Производная & Ошибка в числителе производной, поиск максимума не завершён & Не указан & 562,5 Вт\\
\bottomrule
\end{tabularx}
\normalsize
\vspace{15pt}
\textbf{Как работать с разбором.} Сначала сравните свой последний верный шаг с первым неверным переходом. Затем восстановите решение самостоятельно и выполните одну задачу на закрепление в конце соответствующего разбора.
\clearpage
\thispagestyle{empty}
\vspace*{0.36\textheight}
\begin{center}
{\Huge\color{accent} Разбор ошибок}\par
\vspace{9pt}
{\large\color{quiet} Пять заданий для повторения}
\end{center}
\clearpage
\problem{6}{Математическое ожидание выигрыша}
\blockhead{Условие}
Всего выпущено 50\,000 лотерейных билетов. Из них 1000 выигрывают по 100 руб., 100 билетов --- по 500 руб., 50 билетов --- по 1000 руб., 10 билетов --- по 50\,000 руб. и 2 билета --- по 100\,000 руб. Найдите математическое ожидание выигрыша одного случайно купленного билета.
\blockhead{Ваше решение и ответ}
Записана сумма всех выплат:
\[
100\cdot1000+500\cdot100+1000\cdot50+50\,000\cdot10+100\,000\cdot2=900\,000.
\]
В ответе указано \(900\,000\) руб.
\blockhead{Где именно возникает ошибка}
\err{Верно найдена общая сумма всех денежных выигрышей, однако она названа математическим ожиданием выигрыша одного билета.} Первый неверный шаг --- отождествление суммы \(900\,000\) руб. со средним выигрышем случайного покупателя.
\blockhead{Почему так делать нельзя}
Математическое ожидание --- среднее значение случайной величины с учётом вероятностей. Общая сумма денег относится ко \emph{всем} 50\,000 билетам. При случайном выборе одного билета каждый билет имеет вероятность \(1/50\,000\). При этом большинство билетов невыигрышные; для них выигрыш равен нулю.
\idea{Нужно разделить совокупную сумму выигрышей на общее количество билетов, включая невыигрышные.}
\blockhead{Исправление и полное решение}
Обозначим выигрыш одного билета через \(X\). Тогда
\begin{align*}
M(X)={}&100\cdot\frac{1000}{50\,000}
+500\cdot\frac{100}{50\,000}
+1000\cdot\frac{50}{50\,000}\\
&+50\,000\cdot\frac{10}{50\,000}
+100\,000\cdot\frac{2}{50\,000}
+0\cdot\frac{48\,838}{50\,000}.
\end{align*}
Вынесем общий знаменатель:
\[
M(X)=\frac{100\,000+50\,000+50\,000+500\,000+200\,000}{50\,000}
=\frac{900\,000}{50\,000}=18.
\]
\blockhead{Проверка}
Выигрышных билетов \(1000+100+50+10+2=1162\); невыигрышных --- \(50\,000-1162=48\,838\). Всего действительно 50\,000 билетов. Средняя выплата 18 руб. на один билет означает, что на 50\,000 билетов приходится \(18\cdot50\,000=900\,000\) руб., что согласуется с таблицей.
\result{Математическое ожидание выигрыша одного билета составляет \(18\) рублей.}
\practice{В лотерее выпущено 10\,000 билетов. Из них 60 дают выигрыш 100 руб., 20 --- 500 руб., 5 --- 1000 руб., остальные --- без выигрыша. Найдите математическое ожидание выигрыша одного случайного билета в рублях.}
\clearpage
\problem{10}{Расстояние до горизонта; точность вычислений}
\blockhead{Условие}
Расстояние \(l\) (в километрах) до горизонта вычисляют по формуле \(l=\sqrt{2Rh}\), где \(R=6400\) км, \(h\) --- высота наблюдателя в километрах. При исходном положении горизонт виден на расстоянии 4 км. На сколько метров нужно подняться, чтобы расстояние увеличилось до 28 км?
\blockhead{Ваше решение и ответ}
Для исходной высоты получено \(1{,}25\) м --- это верно. Далее после возведения второго расстояния в квадрат записано \(786=2\cdot6400\cdot h\). В ответе указано \(60{,}15\) м.
\blockhead{Где именно возникает ошибка}
Последний правильный шаг для второго случая:
\[
28=\sqrt{2\cdot6400\cdot h_2}.
\]
Первый неверный шаг:
\[
28^2=786\qquad\text{(ошибка)}.
\]
\err{Квадрат 28 равен 784: \(28\cdot28=784\).} Последующие вычисления и ответ сдвинулись из-за этой арифметической неточности.
\idea{Вычислить две высоты по одной формуле, перевести их в метры и найти разность.}
\blockhead{Полное правильное решение}
Возведём обе части формулы в квадрат:
\[
l^2=2Rh,\qquad h=\frac{l^2}{2R}.
\]
Для \(l_1=4\) км:
\[
h_1=\frac{4^2}{2\cdot6400}=\frac{16}{12\,800}=0{,}00125\ \text{км}=1{,}25\ \text{м}.
\]
Для \(l_2=28\) км:
\[
h_2=\frac{28^2}{2\cdot6400}=\frac{784}{12\,800}=0{,}06125\ \text{км}=61{,}25\ \text{м}.
\]
Нужная прибавка высоты равна
\[
\Delta h=h_2-h_1=61{,}25-1{,}25=60\ \text{м}.
\]
\blockhead{Проверка}
Если человек поднялся на 60 м от исходного положения, высота равна \(61{,}25\) м, или \(0{,}06125\) км. Тогда
\[
\sqrt{2\cdot6400\cdot0{,}06125}=\sqrt{784}=28\ \text{км}.
\]
\result{Подняться нужно на \(60\) метров.}
\practice{Для той же формулы и радиуса Земли \(R=6400\) км расстояние до горизонта увеличили с 6 км до 18 км. На сколько метров повысилась точка наблюдения?}
\clearpage
\problem{11}{Совместная работа двух труб}
\blockhead{Условие}
Каждая из двух одинаковых труб может заполнить бассейн за 11 часов. Через 3 часа после начала работы первой трубы подключили вторую. За какое время после подключения второй трубы бассейн заполнился полностью?
\blockhead{Ваше решение и ответ}
На представленном бланке решение и ответ отсутствуют.
\blockhead{Где именно возникает затруднение}
\err{Решение не начато, поэтому невозможно выделить первый ошибочный арифметический шаг.} Для выполнения задачи необходимо выразить производительность каждой трубы и учесть объём, заполненный до подключения второй.
\idea{Весь бассейн принять за единицу; скорость работы одной трубы составляет одну одиннадцатую бассейна за час.}
\blockhead{Полное правильное решение}
Одна труба за час заполняет
\[
\frac{1}{11}\ \text{бассейна}.
\]
За первые три часа первая труба заполняет
\[
3\cdot\frac{1}{11}=\frac{3}{11}\ \text{бассейна}.
\]
Значит, до полного заполнения осталось
\[
1-\frac{3}{11}=\frac{8}{11}\ \text{бассейна}.
\]
Теперь работают две одинаковые трубы. Вместе за один час они заполняют
\[
\frac{1}{11}+\frac{1}{11}=\frac{2}{11}\ \text{бассейна}.
\]
Пусть \(t\) --- число часов \emph{после подключения второй трубы}. Тогда
\[
\frac{2}{11}t=\frac{8}{11},
\qquad t=\frac{8}{11}:\frac{2}{11}=4.
\]
\blockhead{Проверка}
Первая труба работала \(3+4=7\) часов и заполнила \(7/11\) бассейна; вторая работала 4 часа и заполнила \(4/11\). Их общий результат:
\[
\frac{7}{11}+\frac{4}{11}=1.
\]
\result{После подключения второй трубы прошло \(4\) часа.}
\practice{Каждая из двух одинаковых труб заполняет бассейн за 15 часов. Вторую трубу включили через 3 часа после первой. За сколько часов после подключения второй трубы бассейн заполнится?}
\clearpage
\problem{16}{Логарифмическое неравенство и область допустимых значений}
\blockhead{Условие}
Решить неравенство
\[
\log_{0{,}5}\!\left(\frac{x-3}{x+3}\right)+2\log_4(x^2-9)\ge 5.
\]
\blockhead{Ваше решение и ответ}
Верно выполнен переход к основанию 2 и объединены логарифмы. После получения
\[
\log_2 (x+3)^2\ge5
\]
в решении появляется \((x-3)^2\ge32\), затем записан отрезок
\(\left[3-4\sqrt2;\,3+4\sqrt2\right]\).
\blockhead{Где именно возникает ошибка}
Последний правильный результат: \(\log_2 (x+3)^2\ge5\).
\err{При снятии логарифма \(x+3\) заменено на \(x-3\). Далее неравенство вида «квадрат не меньше положительного числа» решено как промежуток между корнями. Также не записана область допустимых значений.}
\blockhead{Почему это существенно}
Из \(u^2\ge32\) следует \(u\le-4\sqrt2\) или \(u\ge4\sqrt2\): величина по модулю должна быть достаточно большой. Точки \(x=\pm3\) запрещены исходными логарифмами, даже если после сокращения множители исчезли.
\idea{Сначала определить ОДЗ, затем преобразовать логарифмы и пересечь полученное множество с ОДЗ.}
\blockhead{Полное правильное решение}
\textbf{1. Область допустимых значений.} Оба аргумента должны быть положительными:
\[
\frac{x-3}{x+3}>0,\qquad x^2-9>0.
\]
Каждое условие выполняется при \(x<-3\) или \(x>3\). Поэтому
\[
\boxed{\text{ОДЗ}:\quad x\in(-\infty;-3)\cup(3;+\infty).}
\]
\textbf{2. Приведём логарифмы к основанию 2.}
\[
\log_{0{,}5}A=-\log_2 A,
\qquad 2\log_4 B=\log_2 B\quad(B>0).
\]
Следовательно, на ОДЗ:
\begin{align*}
-\log_2\!\left(\frac{x-3}{x+3}\right)+\log_2(x^2-9)&\ge5,\\
\log_2\!\left(\frac{x^2-9}{(x-3)/(x+3)}\right)&\ge5.
\end{align*}
Поскольку \(x\ne3\), можно сократить:
\[
\frac{(x-3)(x+3)}{(x-3)/(x+3)}=(x+3)^2.
\]
Значит, \(\log_2(x+3)^2\ge5\).
\textbf{3. Решим полученное неравенство.} Основание 2 больше единицы, поэтому знак сохраняется:
\[
(x+3)^2\ge32\iff |x+3|\ge4\sqrt2.
\]
Отсюда
\[
x\le-3-4\sqrt2\qquad\text{или}\qquad x\ge-3+4\sqrt2.
\]
\textbf{4. Пересечём с ОДЗ.} Так как \(-3+4\sqrt2\approx2{,}66<3\), вся положительная ветвь ОДЗ \(x>3\) подходит, а на отрицательной ветви остаётся \(x\le-3-4\sqrt2\).
\[
\boxed{x\in(-\infty;-3-4\sqrt2]\cup(3;+\infty).}
\]
\blockhead{Проверка}
При \(x=4\) исходные логарифмы определены и \((4+3)^2=49\ge32\), точка подходит. При \(x=-4\) аргументы определены, но \((-4+3)^2=1<32\), точка не подходит. Точка \(x=3\) исключена, поскольку исходные логарифмы не определены.
\result{Ответ: \(\displaystyle (-\infty;-3-4\sqrt2]\cup(3;+\infty)\).}
\practice{Решите неравенство \(\displaystyle \log_{0{,}5}\!\left(\frac{x-2}{x+2}\right)+2\log_4(x^2-4)\ge4\).}
\clearpage
\problem{17}{Наибольшая мощность; дифференцирование}
\blockhead{Условие}
Электродвижущая сила аккумулятора \(\mathcal E=60\) В, внутреннее сопротивление \(r=1{,}6\) Ом. Сила тока в подключённой зарядной станции: \(I=\mathcal E/(r+R)\), где \(R\ge0\) --- сопротивление станции. Тепловая мощность \(P=I^2R\). Найдите наибольшую возможную мощность в ваттах.
\blockhead{Ваше решение и ответ}
Верно получена функция
\[
P(R)=\frac{3600R}{(R+1{,}6)^2}.
\]
Далее применено правило дифференцирования дроби. При раскрытии скобок и упрощении числителя получено ошибочное выражение \(-2R^2+2R+2{,}56\); составлено уравнение \(R^2-R-1{,}28=0\). Вычисления обрываются на дискриминанте, окончательный ответ не указан.
\blockhead{Где именно возникает ошибка}
\err{В числителе производной неверно раскрыт и упрощён квадрат \((R+1{,}6)^2\).} Именно это приводит к неправильному квадратному уравнению. Кроме того, даже найденную стационарную точку нужно было бы проверить на максимум и вычислить саму мощность.
\idea{Можно использовать производную, а можно сразу оценить знаменатель с помощью квадрата разности. Оба школьных способа приводят к одному максимуму.}
\blockhead{Полное правильное решение: способ с производной}
Подставим выражение для силы тока в мощность:
\[
P(R)=I^2R=\left(\frac{60}{1{,}6+R}\right)^2R
=\frac{3600R}{(1{,}6+R)^2},\qquad R\ge0.
\]
По правилу производной частного:
\begin{align*}
P'(R)
&=\frac{3600(R+1{,}6)^2-3600R\cdot2(R+1{,}6)}{(R+1{,}6)^4}\\
&=\frac{3600(R+1{,}6)\bigl[(R+1{,}6)-2R\bigr]}{(R+1{,}6)^4}\\
&=\frac{3600(1{,}6-R)}{(R+1{,}6)^3}.
\end{align*}
Знаменатель положителен при всех \(R\ge0\). Значит,
\begin{itemize}
\item при \(0\le R<1{,}6\) производная положительна: мощность возрастает;
\item при \(R=1{,}6\) производная равна нулю;
\item при \(R>1{,}6\) производная отрицательна: мощность убывает.
\end{itemize}
Следовательно, \(R=1{,}6\) Ом --- точка максимума:
\[
P_{\max}=P(1{,}6)=\frac{3600\cdot1{,}6}{(1{,}6+1{,}6)^2}
=\frac{5760}{10{,}24}=562{,}5\ \text{Вт}.
\]
\blockhead{Дополнительная проверка без производной}
Из \((R-1{,}6)^2\ge0\) для \(R\ge0\) получаем
\[
(R+1{,}6)^2-6{,}4R=(R-1{,}6)^2\ge0,
\]
то есть \((R+1{,}6)^2\ge6{,}4R\). При \(R>0\)
\[
P(R)=\frac{3600R}{(R+1{,}6)^2}\le\frac{3600R}{6{,}4R}=562{,}5.
\]
Равенство достигается при \(R=1{,}6\). При \(R=0\) мощность равна нулю, так что максимум подтверждён на всей допустимой области.
\result{Наибольшая тепловая мощность равна \(562{,}5\) Вт при \(R=1{,}6\) Ом.}
\practice{Для зарядной станции \(\mathcal E=48\) В, \(r=2\) Ом, \(I=\mathcal E/(r+R)\) и \(P=I^2R\), где \(R\ge0\). Найдите наибольшую тепловую мощность и сопротивление, при котором она достигается.}
\end{document}
